\chapter{Implementation}
\label{ch:implementation}

\section{Environment Setup}
\label{sec:impl_env}

% TODO: Hardware (RTX 5080), OS (Windows 11), Python 3.14, PyTorch 2.11,
% acids-rave 2.3.1, pytorch-lightning 2.6.1, Pure Data with nn~ external.

\subsection{Compatibility Patches}
% TODO: List the four patches applied to acids-rave 2.3.1 to make it run
% with our newer torch/scipy/lightning stack:
%   1. pqmf.py: scipy.signal.kaiser moved to scipy.signal.windows.
%   2. pqmf.py: kaiserord argument must be a Python float.
%   3. pqmf.py: firwin nyq parameter renamed to fs.
%   4. model.py: validation_epoch_end -> on_validation_epoch_end (Lightning 2.0).
% These are reproducible in any fresh install and are documented in
% docs/setup_notes.md.

\section{Training Pipeline}
\label{sec:impl_training}

\subsection{Configuration Files}
% TODO: Describe the .gin configuration system. We started with custom
% c16_r10 (capacity 64, ratios [2,2,2,1], LATENT_SIZE 128 then 16). After
% repeated GAN failures, we pivoted to the official RAVE v2 + causal
% configs with only two overrides (LATENT_SIZE = 16, PHASE_1_DURATION =
% 1.5M). The shift was decisive: v6 produced a working autoencoder.

\subsection{Auto-Resume Training Scripts}
% TODO: All training is launched via Windows .bat scripts that:
%   1. Activate the venv.
%   2. Search for the latest checkpoint via PowerShell.
%   3. Resume from it if found, otherwise start fresh.
%   4. Pass --val_every 10000 for regular checkpointing.
% This makes training safe to interrupt and resume without losing state.

\subsection{Validation Checkpointing}
% TODO: best.ckpt selected by validation reconstruction loss; epoch-tagged
% checkpoints saved every 500k steps. This lets us recover the model from
% any point in training, useful when Phase 2 degrades it later.

\section{Real-Time Deployment}
\label{sec:impl_realtime}

\subsection{Model Export}
% TODO: \texttt{rave export --run RUN\_DIR --name NAME} produces a
% TorchScript .ts file compatible with the nn~ external. Internally,
% export selects best.ckpt (lowest validation loss) by default.

\subsection{Pure Data Patches}
% TODO: Reference figures of the test patches:
%   - test_drums.pd: mic -> nn~ drums_v1 forward -> *~ 50 -> dac~
%   - test_guitar_v6_final.pd: mic -> nn~ guitar_v6 forward -> *~ 30 -> dac~
% Note the gain compensation: BRAVE output amplitude is in [0.005, 0.05],
% so a fixed gain multiplier is needed to reach normal listening levels.

\subsection{Measured Latency}
% TODO: Pure Data block size 64 samples at 44.1 kHz = 1.45 ms per buffer.
% Plus model forward time (typically < 1 ms on a modern GPU). Plus audio
% driver buffer (Windows WASAPI ~5 ms). Total end-to-end < 10 ms,
% confirming the BRAVE design target.

\section{Evaluation Tooling}
\label{sec:impl_eval}

\subsection{Diagnostic Script}
% TODO: scripts/analyze_noise_floor.py and ad-hoc Python tests run in
% the venv. Tests are reproducible and produce both numerical metrics
% and WAV files for qualitative listening.

\subsection{TensorBoard Monitoring}
% TODO: \texttt{tensorboard --logdir runs/} watches all runs. Critical
% scalars are pinned for quick health checks at each training session.
